\paragraph{\trSuspension{}操作与A∞结构。}\label{sub:shift-and-a-infinity-structure}
我们现在要说明，\trExactLagrangianBrane{}上的\trSuspension{}操作——由\eqref{eq:shift-operation} 逐点定义——确实对应于正确的代数概念；更明确地说，它表示相应的那种\trAinfty-模，即在\cref{sub:shift}中作为\eqref{eq:tensor-product-module} 的特例所定义的模。这一事实在\trCohomology{}层面的结论，已在我们讨论\trFloerCohomology{}时陈述过（未加证明），见\eqref{eq:brane-shift-iso} 与\eqref{eq:brane-shift-product-diagram}。
考虑一对\trExactLagrangianBrane{}$L_0^{\#}$、$L_1^{\#}$，以及后者经\trSuspension{}的版本 $SL_1^{\#}$。取某个 $y\in\scr{C}(L_0,L_1)$。一个有用的记号区分是：把同一条路径 $y$ 在联想到\trBrane{}对 $(L_0^{\#},SL_1^{\#})$ 时记作 $\tilde{y}$。这意味着我们同时有\trIndex{}与\trOrientation{}空间 $i(\tilde{y})$、$o(\tilde{y})$ 以及 $i(y)$、$o(y)$。\cref{thm:shift-index-orientation}说，$i(\tilde{y})=i(y)-1$，并且存在一个\trCanonical\trIsomorphism
\begin{equation}\label{eq:shift-orientation-isomorphism}
\sigma_y\colon o(y)\longrightarrow o(\tilde{y}).
\end{equation}
回顾该引理的证明便可记起：在\trIndex{}理论的语言下，\eqref{eq:shift-orientation-isomorphism} 是如下事实的一个推论——\trOrientationOperator{}$D_{H,\tilde{y}}$ 是 $D_{H,y}$ 与\trDisk{}上一个\trIndex{}为 $n-1$ 的算子 $D_{D,\rho}$ 的\trBoundaryConnectedSum{}（依此顺序）。事实上，把 $D_{D,\rho}$ 的\trDomain{}限制到在 $\partial D$ 的某个取定点处为零的函数所成的空间上或许更好，因为这样一来\trIndex{}便降到 $-1$，同时还去掉了\trGluing{}公式\eqref{eq:boundary-gluing} 中额外的 $\lambda^{\mathrm{top}}(F)$ 项。
我们的基本论断是：若 $u\in\scr{M}^{*}_Z(y_0,y_1)$ 是一条\trIsolated{}的\trFloer{}\trTrajectory，而 $\tilde{u}$ 是同一个映射视作 $\scr{M}^{*}_Z(\tilde{y}_0,\tilde{y}_1)$ 中元素时的样子，则与之相伴的\trIsomorphism{} \eqref{eq:relative-orientation-cu} 与\eqref{eq:shift-orientation-isomorphism} 交换，也就是说
\begin{equation}\label{eq:sigma-commutation-claim}
\sigma_{y_0}\circ c_u=c_{\tilde{u}}\circ\sigma_{y_1}.
\end{equation}
两边都是把 $D_{H,\tilde{y}_0}$ 设想成由 $D_{Z,u}$（即 $u$ 的\trLinearization{}算子）、$D_{H,y_1}$ 与 $D_{D,\rho}$ \trGluing{}而成的结果。由于这些因子两次出现的顺序相同，\eqref{eq:sigma-commutation-claim} 可由\eqref{eq:delbar-gluing} 与\eqref{eq:boundary-gluing} 的结合性质推出。我们用\eqref{eq:shift-orientation-isomorphism} 的 $\KK$-\trNormalization{}、乘上 $(-1)^{i(y)}$，来等同
\begin{equation}\label{eq:shifted-cochain-iso}
CF^*(L_0^{\#},SL_1^{\#})\cong CF^{*+1}(L_0^{\#},L_1^{\#}).
\end{equation}
按定义，它与 $\partial$ 反交换，因而与 $\mu^1$ 交换。
现在考虑一个由\trBrane{}组成的 $(d{+}1)$ 元组 $(L_0^{\#},\ldots,L_d^{\#})$（$d\ge 2$），以及一个\trIsolated{}点 $u\in\scr{M}^{d+1}(y_0,\ldots,y_d)$。我们要把 $L_d^{\#}$ \trSuspension{}成 $SL_d^{\#}$，并照前样引入附加记号 $\tilde{y}_0$、$\tilde{y}_d$、$\tilde{u}$。与\eqref{eq:sigma-commutation-claim} 平行的断言说，诸\trOrientation{}空间之间的\trMultilinear{}映射 $c_u$ 与 $c_{\tilde{u}}$ 满足
\begin{equation*}
\sigma_{y_0}\circ c_u=(-1)^{\star}\,c_{\tilde{u}}\circ(\sigma_{y_d}\otimes\id\otimes\cdots\otimes\id),
\end{equation*}
其中 $\star=i(y_1)+\cdots+i(y_{d-1})$。论证与前面相同，只是在联系两边时，须按\trComposition{}的次序把 $D_{D,\rho}$ 移过 $D_{H,y_1},\ldots,D_{H,y_{d-1}}$，这就产生了 Koszul 型的符号。于是，$\mu^d$ 与\eqref{eq:shifted-cochain-iso} 交换——这恰恰就是一般\trACategory{}中\trSuspension{}后\trObject{}的定义。为看清这一点，注意除 $\star$ 之外，其他有关的符号量还有：由 $y_d$ 与 $y_0$ 的\eqref{eq:shifted-cochain-iso} 定义而来的 $i(y_d)$ 与 $i(y_0)=i(y_1)+\cdots+i(y_d)+2-d$；此外还有由\eqref{eq:dagger-sign-exponent} 而来的 $i(y_1)+\cdots+d\,i(y_d)$ 与 $i(y_1)+\cdots+d\,i(\tilde{y}_d)$。最终所有符号都互相抵消，得到想要的结果。
